\begin{frame}{확인 문제 (15.1)}
  \begin{enumerate}
    \item \textbf{ELBO 부등식의 갭.} $\log p(x) - \text{ELBO} =
      D_{KL}(q \| p(z|x))$임을 유도한 뒤, $q = p(z|x)$일 때만 갭이 0이
      된다는 것을 한 문장으로 설명.
    \item \textbf{정규분포 합 = 정규분포.}
      $\mathcal{N}(x; z,1) \cdot \mathcal{N}(z;0,1)$에 대해 $z$로 적분
      하면 $\mathcal{N}(x; 0, 2)$가 된다는 것을, ``두 독립 정규분포의
      합은 정규분포''라는 사실로부터 한 문장으로 설명.
    \item \textbf{손계산 재현.} $x = 9$ 예제에서 복원 항이
      $\approx -33.4$가 되는 것을 직접 계산 (힌트: $(x-\mu)^2 = 64$,
      $\sigma_q^2 = 1$, $\log 2\pi \approx 1.838$ 대입).
  \end{enumerate}
  \vskip0.8em
  \begin{block}{핵심 요약}
    ELBO = \textbf{복원 항} $-$ \textbf{KL 정규화 항}.
    \textbf{ELBO 최대화 = 사후분포 근사 $q$를 정확히 만드는 것
    (변분 베이즈)}. \; $\mu^2$ 항은 \textbf{위치를 원점 쪽으로
    끌어당기는 힘}, $\sigma^2$ 항은 \textbf{너비를 제어}.
  \end{block}
\end{frame}
